\scanpage{165}
\mspage{477}
% position: note 1082 upper paragraph
\(({}^{108}2)\) (6 octobre) Dans ce ``déployement de force'' il n'y a nulle intention
``agressive'' au sens courant du terme, consciente ni inconsciente, seulement un
désir inconscient d'impressionner, de forcer l'estime. Il est vrai que ce terme
``forcer l'estime'' qui me revient spontanément, porte déjà une connotation de
\mbox{c o n t r a i n t e}, voisine de celle ``d'agression''. Cette intention inconscien-\linebreak
te de contrainte, perçue au niveau inconscient également, doit souvent être \unsure{vécue}
comme une sorte d'agression (alors même que ce vécu reste occulté, tout comme \unsure{les}
réactions antagonistes qu'il déclenche). En même temps, il doit se faire souvent
l'amalgame de ce vécu avec des vécus analogues, remontant à l'enfance, avec le
père comme protagoniste, et où celui-ci apparaît comme le principal détenteur
d'autorité répressive, voire comme un rival écrasant, envié et détesté.

% position: middle paragraph
Même sans un tel amalgame, et indépendamment aussi de toute perception \unsure{chez}
autrui d'une intention de ``contrainte'' en moi, il doit y avoir souvent la per-\linebreak
ception d'un fort \mbox{d é s é q u i l i b r e}, d'une disharmonie foncière, dans
ce ``déployement de force'' exclusivement yang (dans son esprit et intention, au
moins). Cette démesure est néfaste au principal intéressé, savoir moi-même, et
à la limite bel et bien ``dangereux'' pour sa survie physique même (comme des inci-\linebreak
dents-santé de ces dernières années me l'ont bien montré! ). C'est cela sans dou-\linebreak
te qui a été en filigrane dans ma pensée, quand j'écrivais qu' ``un tel \unsure{déployemer|déployement}
de force'' était ressenti ``en tout cas comme dangereux'' - dangereux ``par nature'',
un exemple donc à surtout ne pas suivre...! Un tel ressenti est sûrement suffi-\linebreak
sant pour susciter des ``réactions de défense'', en l'absence même de toute agres-\linebreak
sion ou intention d'agresser.

% position: lower paragraph
Il est vrai que de telles relations d'ambiguité se sont reproduites après
1976, avec certains de mes élèves notamment, en des moments où tout investisse-\linebreak
ment mathématique était absent, et où il n'y avait aucun ``dépoyement de force''
apparent dans ma vie. Il est vrai aussi que les `` déployements'' en question du
\mbox{p a s s é} ont créé une \mbox{r é p u t a t i o n}, qui continue à me coller à la
peau, surtout dans ma vie professionelle, et qui dans une certaine mesure se\inserted{\unsure{,}}
substitue à la perception de celui que je suis \mbox{d a n s\quad l e\quad p r é s e n t}.
De plus, j'ai acquis dans le commerce de certains thèmes mathématiques une aisan-\linebreak
ce telle que, même en dehors de mes périodes mathématiques et ma réputation ai-\linebreak
dant, cette aisance ou maîtrise naturelle peuvent avoir déjà l'effet du ``déploye-\linebreak
ment de force'', sur des élèves peu motivés, et me faire ressentir par eux (en
dépit de certains traits avenants voire rassurants) comme une sorte de Superman
(un peu Superpère sur les bords!).
% position: left lower mark
\margin{\unsure{-}}
% position: right foot pencil rotated
\margin{164}
